\section{Discussion}
\label{sec:discussion}

We have shown that a per-class conformal correction can be predicted from label-free
descriptors of the class and applied to classes that have no calibration examples at all. On
ImageNet it takes $70\%$ of an oracle ceiling on those classes, improves on every existing
method we were able to run, two of which are undefined for all of them, and needs labels for as
few as a tenth of the classes. Marginal coverage is preserved because only a single scalar is
calibrated (\cref{prop:marginal}).

\paragraph{Limitations.} There is no finite-sample class-conditional guarantee at $n_y=0$, for
our method or for any other, and our claim for those classes is therefore empirical. The
improvement attenuates for adaptive score functions and disappears at $\alpha=0.01$. We have
also not shown that the approach transfers across architectures: the additional backbones leave
only $35$ evaluation examples per class, which \cref{fig:evaldepth} places below the depth at
which any effect is detectable, so those experiments do not settle the question in either
direction. Settling that needs an evaluation set larger than ImageNet's validation split
provides.

Three further limitations are worth naming precisely. The percentages of the oracle ceiling we
quote are ratios of averages, and the split-level spread behind them is wide: on the primary
arm the fraction of the ceiling taken ranges from $2\%$ to $75\%$ across five splits, because
the ceiling itself moves by $0.10$ between splits. Five splits is what ImageNet's validation
set supports at this evaluation depth, and it is not many. Second, the depth at which the
ceiling collapses was measured on one dataset at one miscoverage level; we validated the
prediction it makes on Pl@ntNet-300K, but we did not check whether the threshold sits at $35$
examples per class for other $\alpha$ or other label spaces, and we would not expect a single
number to transfer exactly. Third, the configuration behind every table lets $g_\theta$,
$\lambda$ and $c$ share calibration rows, which is not what \cref{prop:marginal} assumes; the
measured cost is $0.0024$ of marginal coverage (\cref{app:invisible}) and the variant that
satisfies the premise is reported separately (\cref{app:split}). Fourth, all of our tables use
a random per-class subsample of the score dump, adopted for memory reasons and described in
\cref{app:rowbudget}; it leaves $2.7$ times the depth the pre-registered premise required and,
by both of our own depth results, should if anything understate the effect, but the budget is
swept there so that the claim rests on a measurement.

\paragraph{Hypotheses we discarded.} Three explanations we held during the work were refuted by
experiments designed to test them, and we record them with the measurements that ruled them
out. We first attributed the failures to calibration depth, which its own sweep refuted, since
ten calibration examples per class still yields a significant improvement. We then attributed
the backbone failures to descriptors that were endogenous to the scoring model, which running
both descriptor arms refuted, with every interval containing zero. Finally we attributed the
null result on iNaturalist-2018 to scarce calibration data, whereas the surviving classes have
$67.5$ examples each and the mechanism is instead that shrinkage selects $\lambda=0$ and the
method declines to act.

\paragraph{Evaluating class-conditional coverage.} The observation we expect to be most useful
beyond this particular method is that worst-class coverage becomes unmeasurable before it
becomes unachievable. Below roughly $35$ evaluation examples per class the oracle ceiling is
zero, so a null result on such a slice says little about the method that produced it.
Long-tailed benchmarks sit squarely in that regime, which means published comparisons on them
may be reporting a property of the evaluation set where they intend to report a difference
between methods. We
state the measurability regime alongside every number for this reason, and the practice seems
worth adopting more widely.

\paragraph{Future work.} The most direct next step is a larger evaluation set for the rare
classes themselves, since for Pl@ntNet-300K and iNaturalist-2018 the restriction used in
\cref{tab:longtail} retains under $10\%$ of classes and the interesting question is what happens
to the rest. A second question is why $\varphi$ carries signal at $K=1081$ classes but not at
$K=8142$; the diagnostic in \cref{app:longtail} locates the failure in the regression, since the
shrinkage curve there is informative and points the wrong way, so the selection step is doing
its job. A third
is whether richer descriptors, such as taxonomies, text embeddings of class names, or features of
the training images of related classes, extend the reach of the method, since nothing in
\cref{prop:marginal} requires $\varphi$ to be geometric.

The most interesting direction is a formal statement to replace the empirical one. No
procedure can bound coverage for an individual class with $n_y=0$. Group-conditional
conformal prediction can bound coverage for a \emph{neighbourhood} of classes, however, and the
framework of \citet{bairaktari2025kandinsky} accepts overlapping and fractional groups defined
jointly on covariates and labels, which is exactly the shape our descriptors have. Groups drawn
in $\varphi$-space would give a guarantee about regions of the classifier's output geometry, which
is weaker than a per-class statement and much stronger than none.
